\clearpage
\par\bigskip\pdfbookmark[0]{On the Connection. Fragment III}{nexus-fragment-III}\centerline{[III.]}\medskip
Let $C$ be the collection of representatives of all classes of properly primitive forms for the determinant $-D$. We denote by $(n)$ the number of all representations of the number $n$ by forms in the collection $C$. Let $p$ be an odd prime. Then

\SourceLeaf{293}
\[
\begin{aligned}
&1.\quad(pn)=(n)\quad\text{if }p\text{ divides }D\\
&\left.\begin{aligned}2.\quad(pn)&=(n)+(h)\\3.\quad(pn)&=-(n)+(h)\end{aligned}\right\}
\ \text{if }p\text{ does not divide }D\\
&\hspace{7em}\left\{\begin{aligned}&\text{divides }xx+D\\&\text{does not divide }xx+D\end{aligned}\right.
\end{aligned}
\]
where $n=hp^\mu$, $\mu$ is arbitrary, and $h$ is not divisible by $p$.
\[
\begin{aligned}
\text{In case }1.\ \text{we have}\quad&(h)=(ph)=(pph)=(p^3h)\text{ etc.}\\
2.\quad&(ph)=2(h),\ (pph)=3(h),\ (p^3h)=4(h)\text{ etc.}\\
3.\quad&(ph)=0,\ (pph)=(h),\ (p^3h)=0,\ (p^4h)=(h)\text{ etc.}
\end{aligned}
\]
\begin{center}$\ast\qquad\ast$\par$\ast$\end{center}

From each properly primitive class for the determinant $=-D$, the number of these classes being $=\lambda$, choose one form, and let $L$ be the collection of these forms.

Let $fA$ denote the number of all representations of the number $A$ by forms in $L$.

Further, let $f(A;p)=f\frac A{p^\alpha}$ when $p^\alpha$ is the highest power of the prime $p$ that divides $A$; also $f(A;p,q)=f\frac A{p^\alpha q^\beta}$ when $q$ is another prime whose highest power dividing $A$ is $=q^\beta$; and similarly $f(A;p,q,r)=f\frac A{p^\alpha q^\beta r^\gamma}$ when $r$ is a third prime whose highest power dividing $A$ is $r^\gamma$, and so on.
\BookRule{.30\textwidth}
